% Abhakseite „Das kann ich“ – Zone nur im Gesamt (\mitzone)
%% TODO Ich-kann-Sätze: Platzhalter wie in den Aufgabendateien
\begin{abhakseite}
\ifdefined\mitzone
\abhakgruppe{Kennst du schon}
\abhak{1}{Potenzen mit natürlichem Exponenten lesen und schreiben, Vorfaktor und Exponent unterscheiden, Potenzgesetze (Produkt gleicher Basen, Potenz einer Potenz), negative Exponenten als Kehrwert, Wurzel als Potenz mit Exponent ein Halb}
\abhak{2}{Terme umformen}
\abhak{3}{Bruch- und Dezimalkoeffizienten multiplizieren (ein Viertel mal fünf, null Komma vier mal vier), mit Brüchen kürzen}
\abhak{4}{Die e-Funktion als Funktion, die sich selbst als Ableitung hat; $e^{(kx)}$ als Verkettung lesen; $e^x$ ist nie null und immer positiv}
\abhak{5}{Ableitung als Tangentensteigung und Änderungsrate deuten, Ableitungswert an einer Stelle als Steigung lesen, Extrempunkt als Stelle mit waagerechter Tangente}
\abhak{6}{Funktionswerte berechnen und Punktprobe, Anstieg an einer Stelle als Wert der ersten Ableitung}
\abhak{7}{Verschiebung und Streckung eines Graphen an der Gleichung erkennen (f(x − c) verschiebt nach rechts, c · f(x) streckt)}
\abhak{8}{Potenzen mit natürlichem Exponenten lesen und schreiben, Vorfaktor und Exponent unterscheiden, Potenzgesetze (Produkt gleicher Basen, Potenz einer Potenz), negative Exponenten als Kehrwert, Wurzel als Potenz mit Exponent ein Halb – Fehler finden}
\abhak{9}{Potenzen mit natürlichem Exponenten lesen und schreiben, Vorfaktor und Exponent unterscheiden, Potenzgesetze (Produkt gleicher Basen, Potenz einer Potenz), negative Exponenten als Kehrwert, Wurzel als Potenz mit Exponent ein Halb – selbst rechnen}
\fi
\abhakgruppe{Einheit 1 · Potenz-, Faktor- und Summenregel}
\abhak{10}{Zahl oder Variable?}
\abhak{11}{Potenz-, Faktor- und Summenregel}
\abhak{12}{Verhalten im Unendlichen bestimmen}
\abhak{13}{Potenz-, Faktor- und Summenregel – Fehler finden, Begründen}
\abhakgruppe{Einheit 2 · Kettenregel}
\abhak{14}{Kettenregel}
\abhak{15}{Term einer höheren Ableitung einer Sinusfunktion über die Periodizität der Ableitungen angeben}
\abhak{16}{Kettenregel – Fehler finden, Begründen}
\abhakgruppe{Einheit 3 · Produktregel}
\abhak{17}{Produktregel}
\abhak{18}{Produktregel – Fehler finden, Begründen}
\end{abhakseite}
